\section{Technical Appendix Scope}

This appendix records the deterministic HJB proof details that are too long
for the main article.  It is not a second list of assumptions.  The singular
billiard estimates used below are exactly the response-theory outputs imported
in the main paper: reduced observable resolvents, moving-boundary trace
insertions, uniform differentiated resolvent words, suspension contour bounds,
and finite regularity loss.  The present appendix checks how those imported
objects enter the finite-response action, the corrector hierarchy, the
contact-frozen perturbed test, the half-relaxed viscosity passage, and the
nonconvexity algebra.

\begin{definition}[HJB appendix dependency order]
\label{def:supp_hjb_dependency_order}
The proof details below use the following order.
\begin{enumerate}[label=\textup{(\roman*)},wide,labelindent=0pt]
  \item The compact finite-response port fixes the table selector
  \(\Theta(x,p)\), the centered drift \(B(x,p,z)\), and the action read-out
  \(V(x,p,z)\) before any effective coefficient is defined.
  \item Paper 1 supplies the invariant projector, reduced flow resolvent, and
  moving-response bounds for the selected billiard fiber.
  \item The first and second cell equations are solved in the zero-mean
  observable space of that fiber.
  \item The prelimit Hamilton-Jacobi operator is tested only after the actual
  viscosity contact jet has been frozen.
  \item The effective HJB equation is identified by the invariant part of the
  order-one residual; the centered part is cancelled by the second corrector.
  \item The upper and lower relaxed limits are compared by the macroscopic
  viscosity comparison theorem.
\end{enumerate}
Thus the HJB equation is an output of the deterministic calculation, not an
input to the microscopic port or to the response theorem.
\end{definition}

\section{Finite-Response Micro-Action}

\begin{definition}[Microscopic finite-response port]
\label{def:supp_micro_action_port}
The primitive finite-response device consists of a compact internal energy
surface \(\Ycal\), a deterministic finite-time response flow
\[
  \Rcal_s^{x,P}:\Ycal\to\Ycal,\qquad 0\le s\le s_{\rm fr},
\]
a smooth saturation \(\Pi_{\rm fr}:\R^m\to\R^m\), and branchwise smooth read-out
maps
\[
  \vartheta(x,y),\qquad
  \beta(x,y,z),\qquad
  \Vcal(x,y,z).
\]
With \(y_0\) the calibrated rest state and
\(y(x,P)=\Rcal_{s_{\rm fr}}^{x,P}y_0\), set
\[
  \Theta(x,p)=\vartheta(x,y(x,\Pi_{\rm fr}(p))),\quad
  \widehat B(x,p,z)=\beta(x,y(x,\Pi_{\rm fr}(p)),z),\quad
  V(x,p,z)=\Vcal(x,y(x,\Pi_{\rm fr}(p)),z).
\]
The centered drift used in the moving slow frame is
\[
  B(x,p,z)=\widehat B(x,p,z)
   -\int \widehat B(x,p,\zeta)\,d\mu_{\Theta(x,p)}(\zeta).
\]
All derivatives needed below are bounded on compact contact-jet windows and in
homogeneous flow boxes.  On the physical cotangent window
\(\Pi_{\rm fr}(p)=p\), so local Taylor coefficients in \(p\) are genuine
mechanical response coefficients.
\end{definition}

\begin{proposition}[No effective-coefficient feedback]
\label{prop:supp_no_effective_feedback}
The data \(\Theta,B,V\) in \cref{def:supp_micro_action_port} are fixed before
the correctors, Green-Kubo tensor, Hamiltonian, HJB solution, and
theta-expectation semigroup are introduced.  The deterministic construction
factors as
\[
  (\Ycal,\Rcal,\Pi_{\rm fr},\vartheta,\beta,\Vcal)
  \longmapsto
  (\Theta,B,V)
  \longmapsto
  (\Pi_\theta,\Rcal_\theta,\chi,\psi)
  \longmapsto
  (D,H),
\]
where \(\Rcal_\theta\) denotes the reduced observable resolvent supplied by
Paper 1.  No arrow points from \(D,H\), from a viscosity solution \(u\), or from
\(\Etheta\) back to the finite-response port.
\end{proposition}

\begin{proof}
The first arrow is a finite-dimensional deterministic operation: solve the
internal response flow for the finite time \(s_{\rm fr}\), read the table
parameter, force, and action density, and subtract the instantaneous
microcanonical velocity.  The second arrow is the companion response theorem:
for the table \(\Theta(x,p)\), it constructs the invariant projector and the
zero-mean reduced resolvent.  The third arrow solves the cell equations and
takes invariant pairings.  Since the HJB operator appears only in the last
arrow, it cannot be used to define any earlier object.
\end{proof}

\begin{lemma}[Exact endpoint jets from a compact port]
\label{lem:supp_exact_endpoint_jets}
Fix a compact physical cotangent window \(K_{\rm phys}\) and a finite order
\(J\le r-6\).  Let \(\mathcal J(P)\) be any finite list of saturated polynomial
jets in \(P\) satisfying the microscopic parity, zero-gauge, and smallness
constraints of the finite-response port.  Then the compact internal response
system can be enlarged by finitely many coordinates so that its endpoint map
\[
  E_A(x,P)=y(s_{\rm fr};x,P,A)
\]
has exactly the prescribed \(J\)-jet on \(K_{\rm phys}\).  The construction is
triangular in the jet degree and is completed before invariant densities,
cell equations, or effective coefficients are introduced.
\end{lemma}

\begin{proof}
Choose saturated monomials \(M_\alpha(P)=\Pi_{\rm fr}(P)^\alpha\), which agree
with ordinary monomials on \(K_{\rm phys}\).  In a compact coordinate ball of
the internal energy surface, add linear forcing coordinates with amplitudes
\(A_{\alpha\ell}(x)M_\alpha(P)\).  Picard iteration of the deterministic
finite-time ODE gives
\[
  E_A(x,P)=\frac{s_{\rm fr}^2}{2}
  \sum_{\alpha,\ell}A_{\alpha\ell}(x)M_\alpha(P)e_\ell+O(|A|^2).
\]
At degree \(q\) the lower-degree amplitudes have already determined a finite
polynomial error in the endpoint jet.  Adding a fresh internal coordinate with
linear force equal to the negative of that error changes no lower-degree jet
and cancels the degree-\(q\) error.  The finite Jacobian from the new
amplitudes to the degree-\(q\) endpoint coefficients is diagonal with
nonzero entries \(s_{\rm fr}^2/2\), so the recursion terminates.

The local potential is then cut off inside the internal compact energy surface
and saturated outside \(K_{\rm phys}\).  Smallness of the endpoint amplitudes
keeps the scatterer deformation in the geometric window.  The construction is
therefore a finite-dimensional endpoint controllability statement; it does not
use the transfer spectrum, the Green-Kubo tensor, the Hamiltonian, or a
limiting HJB solution.
\end{proof}

\begin{proposition}[Endpoint reciprocity ledger]
\label{prop:supp_endpoint_reciprocity}
The table displacement \(\Theta\), centered drift \(B\), and action read-out
\(V\) are generated by one finite-response action.  On each homogeneous branch
there are zero-mean branch gauges \(R_{\rm fr}^0,R_{\rm fr}^a\) such that
\[
  B^a(x,p,z)
  =
  -\partial_{x_a}\Sfr(x,\Pi_{\rm fr}(p);z)
  +\Acal_{\Theta(x,p)}R_{\rm fr}^a(x,p,z),
\]
\[
  V(x,p,z)
  =
  \partial_{s_{\rm fr}}\Sfr(x,\Pi_{\rm fr}(p);z)
  +\Acal_{\Theta(x,p)}R_{\rm fr}^0(x,p,z),
\]
and
\[
  D_p\Theta(x,p)
  =
  D_y\vartheta(x,E(x,\Pi_{\rm fr}(p)))
  D_PE(x,\Pi_{\rm fr}(p))D\Pi_{\rm fr}(p).
\]
Consequently the finite Taylor coefficients used by the concrete port obey
ordinary action reciprocity identities; they are not independent effective HJB
coefficients.
\end{proposition}

\begin{proof}
Hamilton's endpoint variation formula for the finite-response action gives the
slow force from \(x\)-variation, the running action from response-time
variation, and the table displacement from endpoint variation.  Billiard
flight endpoint changes add only generator coboundaries, represented by the
gauge terms above; their invariant means vanish because
\(\int\Acal_\theta R\,d\mu_\theta=0\) on the generator domain.  Differentiating
the endpoint formulas gives the stated reciprocity relations.  Since
\cref{lem:supp_exact_endpoint_jets} realizes the endpoint jet before any
cell-resolvent object is formed, these relations are microscopic constraints
rather than post-hoc choices of \(D\) or \(H\).
\end{proof}

\begin{definition}[Prelimit viscosity-duality operator]
\label{def:supp_prelimit_operator}
On a fixed contact-jet window, the branchwise prelimit operator has the fast
part
\[
  \Lcal^\varepsilon_{\rm fast}\Phi
  =
  -\varepsilon^{-2}\Acal_{\Theta(x,p)}\Phi
  -\varepsilon^{-1}B(x,p,z)\cdot D_x\Phi ,
\]
and an order-one action part
\[
  \Lcal^\varepsilon_{\rm act}\Phi
  =
  -V^\circ(x,p,z)
  +\text{finite-response work terms generated by the full slow gradient}.
\]
Here \(p\) is the physical cotangent signal read by the port.  In the viscosity
argument \(p\) is frozen at the actual contact gradient before the singular
fast terms are cancelled.  The operator is interpreted in anisotropic duality
against admissible densities in the billiard fiber.
\end{definition}

\section{Exact Prelimit Action Graph}

\begin{definition}[Branchwise action graph]
\label{def:supp_branchwise_action_graph}
Fix \(\varepsilon>0\), a terminal payoff \(\phi\), and a regular homogeneous
flow box in the universal billiard bundle.  A prelimit action graph is a
branchwise \(C^1\) graph
\[
  \Lambda_s^\varepsilon
  =
  \{(x_s,z_s,P_s,A_s):P_s=D_xA_s,\ (x_s,z_s)\in U_{\rm br}\}
\]
whose dynamics on the branch are
\[
  \dot z_s=\varepsilon^{-2}F_{\Theta(x_s,P_s)}(z_s),\qquad
  \dot x_s=\varepsilon^{-1}B(x_s,P_s,z_s),\qquad
  \dot A_s=V^\circ(x_s,P_s,z_s).
\]
The signal \(P_s\) is the cotangent coordinate of this same graph.  It is not
chosen from the later macroscopic solution.
\end{definition}

\begin{proposition}[Closed prelimit action graph]
\label{prop:supp_closed_action_graph}
For every compact contact-jet window and every terminal graph
\(A_T(x,z)=\phi(x)\), \(P_T=D\phi(x)\), the branchwise action graph in
\cref{def:supp_branchwise_action_graph} exists and is unique up to the first
exit time from a homogeneous flow box.  On finite time intervals the graph is
obtained by concatenating finitely many such branch solutions and by ignoring
only the canonical zero-flux singular set.
\end{proposition}

\begin{proof}
On a homogeneous flow box the finite-response read-outs and the pulled-back
billiard vector field are \(C^{r-4}\) in \((x,P,z)\) and uniformly Lipschitz on
the compact window.  The branch ODE therefore has a unique solution.  The
cotangent relation is preserved because the same deterministic action generates
the slow force, the finite-response table displacement, and the running cost:
differentiating the action identity with respect to the initial slow coordinate
gives the adjoint equation for \(P_s\).  Specular reflection preserves the
canonical one-form, since endpoint variations are tangent to the boundary while
the reflected normal velocity changes sign.  Grazing and multiple-collision
hypersurfaces are not places where the pointwise equation is evaluated; they
have zero flux measure and later enter only through weak trace estimates.
\end{proof}

\begin{theorem}[Exact branchwise prelimit Hamilton-Jacobi identity]
\label{thm:supp_exact_prelimit_hj}
On every regular branch, the principal function generated by
\cref{prop:supp_closed_action_graph} satisfies
\[
  -\partial_t U^\varepsilon
  -\varepsilon^{-2}\Acal_{\Theta(x,p)}U^\varepsilon
  -\varepsilon^{-1}B(x,p,z)\cdot p
  -V^\circ(x,p,z)=0,
  \qquad p=D_xU^\varepsilon(t,x,z).
\]
After completion in the observable-duality domain this is the prelimit
viscosity-duality equation used in the main paper.  It is constructed before
the cell problem, the Green-Kubo tensor, the effective Hamiltonian, or the
theta-expectation semigroup.
\end{theorem}

\begin{proof}
For a small branch time \(h>0\), the action identity gives
\[
  U^\varepsilon(t,x,z)
  =
  U^\varepsilon(t+h,x_h,z_h)
  +\int_t^{t+h} V^\circ(x_s,P_s,z_s)\,ds+o(h).
\]
Taylor expansion of the right-hand side along the branch trajectory produces
the time derivative, the fast generator term, and the slow work term
\(\varepsilon^{-1}B\cdot D_xU^\varepsilon\).  Dividing by \(h\) and sending
\(h\downarrow0\) gives the displayed identity.  The anisotropic completion is
obtained by applying this branch identity against admissible densities and
closing the Green identity with the same moving-boundary trace estimates
imported from Paper 1.
\end{proof}

\begin{lemma}[Universal-bundle meaning of a nonconstant contact covector]
\label{lem:supp_universal_bundle_covector}
If a perturbed test has contact covector
\[
  P_\Phi(t,x,z)=D_x\Phi(t,x,z),
\]
then the prelimit operator evaluated at \(P_\Phi\) is a single branch vector
field on the universal billiard bundle.  It is not a simultaneous evolution of
many unrelated billiard tables.  The map
\[
  (x,P,z)\mapsto
  \Acal_{\Theta(x,P)}
  +\varepsilon^{-1}B(x,P,z)\cdot P+V^\circ(x,P,z)
\]
is evaluated along the graph \(z\mapsto P_\Phi(t,x,z)\), and its derivative in
that graph direction is a strong-to-weak moving-response term.
\end{lemma}

\begin{proof}
Normal-coordinate transport identifies nearby tables with one reference
bundle.  In that bundle, the impact equation, reflection law, roof, and
Jacobian are branchwise functions of \((x,P,z)\).  Substituting a branchwise
field \(P_\Phi(z)\) is therefore ordinary composition before the singular
sets are completed.  Differentiating the composed vector field moves the
branch maps and the singular hypersurfaces; these contributions are exactly
the interior, roof, and trace letters controlled by the companion
response theorem.
\end{proof}

\begin{lemma}[Specular endpoint variation]
\label{lem:supp_specular_endpoint_variation}
On a regular homogeneous branch, differentiating the Hamilton principal action
through a collision creates no extra reflection work.  If \(q\) is the
collision point, \(v^-\) and \(v^+\) are the incoming and outgoing unit
velocities, and \(\delta q\in T_q\partial Q_{\Theta(x,p)}\), then
\[
  \langle v^- - v^+,\delta q\rangle=0 .
\]
Thus a regular specular reflection preserves the canonical endpoint term in
the branch action.
\end{lemma}

\begin{proof}
Specular reflection preserves the tangential velocity component and reverses
only the normal component:
\[
  v^+=v^- -2\langle v^-,n(q)\rangle n(q).
\]
The jump \(v^- -v^+\) is therefore normal to the boundary, while a regular
endpoint variation \(\delta q\) is tangent.  Their pairing vanishes.  The
Liouville flux on the two physical traces is the same with opposite normal
orientation, so the Green identity has no physical-boundary residue.
\end{proof}

\begin{proposition}[Moving branch-boundary ledger]
\label{prop:supp_moving_branch_boundary_ledger}
The only boundary terms left after \cref{lem:supp_specular_endpoint_variation}
are derivatives of branch selectors: grazing faces, homogeneity boundaries,
and branch-change hypersurfaces.  In the viscosity-duality pairing these terms
are weak trace letters.  On a compact contact-jet window,
\[
  \sum_{\Gamma\subset\Ssing}
  |\langle [F]_\Gamma,(X_{x,P}\cdot n_\Gamma)\rho\rangle_\Gamma|
  \le
  C_{\rm tr}(M,K)\|F\|_{\#}\|\rho\|_{\Adm_M}
  +C_{\rm tail}(M,K,N),
\]
where \(C_{\rm tail}(M,K,N)\to0\) as the homogeneity cutoff \(N\to\infty\).
\end{proposition}

\begin{proof}
After normal-coordinate trivialization, a moving homogeneous branch is selected
by a characteristic function.  Differentiating it gives a hypersurface
distribution supported on the moved singular curve.  The normal speed of this
hypersurface is controlled by the same impact equation and crossing denominator
used in the response theorem.  Central strips give the trace constant
\(C_{\rm tr}\).  In strip \(m\), the denominator growth is polynomial in \(m\),
while the admissible density tail and the response strip weight are summable.
This gives the cutoff tail \(C_{\rm tail}\).  Hence moving boundaries enter
only through controlled weak traces, not through unaccounted endpoint work.
\end{proof}

\section{Corrector Hierarchy}

\begin{definition}[First corrector and Green-Kubo tensor]
\label{def:supp_first_corrector_gk}
Fix \((x,p)\) and write \(\theta=\Theta(x,p)\).  For the centered drift
components \(B_i(x,p,\cdot)\), define
\[
  \chi_i(x,p,\cdot)
  =
  -\Acal_\theta^{-1}B_i(x,p,\cdot),
  \qquad
  \Pi_\theta\chi_i=0.
\]
For \(a\in\R^m\), the quadratic Green-Kubo form is
\[
  a^T D(x,p)a
  =
  \int_0^\infty
  \int (a\cdot B)(z)(a\cdot B)(\Phi_\theta^t z)
  \,d\mu_\theta(z)\,dt,
\]
equivalently
\[
  D_{ij}(x,p)
  =
  \frac12\int (B_i\chi_j+B_j\chi_i)\,d\mu_\theta
\]
after the finite-response work connection is reduced by the single-action
reciprocity identity.
\end{definition}

\begin{lemma}[First-order cancellation at the actual contact jet]
\label{lem:supp_first_order_cancellation}
Let \(\varphi\in C^3\) and let \((t_\varepsilon,x_\varepsilon)\) be the
prelimit contact point.  Set
\[
  p_\varepsilon=D\varphi(t_\varepsilon,x_\varepsilon),\qquad
  \theta_\varepsilon=\Theta(x_\varepsilon,p_\varepsilon),
\]
and
\[
  \chi_{p_\varepsilon}
  =
  \sum_i p_{\varepsilon,i}
  \chi_i(x_\varepsilon,p_\varepsilon,\cdot).
\]
Then the scale-\(\varepsilon^{-1}\) part of
\(\Lcal^\varepsilon(\varphi+\varepsilon\chi_{p_\varepsilon})\) vanishes:
\[
  -\varepsilon^{-2}\Acal_{\theta_\varepsilon}
  (\varepsilon\chi_{p_\varepsilon})
  -\varepsilon^{-1}B(x_\varepsilon,p_\varepsilon,z)\cdot p_\varepsilon
  =0.
\]
No term of the form
\(\varepsilon^{-1}(p_\varepsilon-p_0)\) is produced, because the cell equation
is solved after the actual prelimit contact jet is known.
\end{lemma}

\begin{proof}
The first corrector solves
\[
  \Acal_{\theta_\varepsilon}\chi_{p_\varepsilon}
  =
  -B(x_\varepsilon,p_\varepsilon,\cdot)\cdot p_\varepsilon .
\]
Substitution gives the displayed cancellation in the observable dual space.
The limiting jet \((x_0,p_0)\) is not inserted until after this cancellation
has been made, so there is no singular mismatch.
\end{proof}

\begin{definition}[Contact-frozen perturbed test]
\label{def:supp_contact_frozen_test}
At the same contact point define
\[
  X_\varepsilon=D^2\varphi(t_\varepsilon,x_\varepsilon),\qquad
  \boldsymbol\chi_\varepsilon
  =
  (\chi_1,\ldots,\chi_m)^T(x_\varepsilon,p_\varepsilon,\cdot).
\]
The first corrector is frozen in its coefficient arguments but multiplied by
the current macroscopic gradient:
\[
  \Xi_\varepsilon[\varphi](t,x,z)
  =
  \sum_i \partial_i\varphi(t,x)
  \chi_i(x,p_\varepsilon,z).
\]
At the contact point its full slow gradient is
\[
  q_\varepsilon(z)
  =
  D_x^{\rm fr}\chi_{p_\varepsilon}(x_\varepsilon,p_\varepsilon,z)
  +X_\varepsilon^T\boldsymbol\chi_\varepsilon(z).
\]
The second term is the Hessian contribution from the complete chain rule; it
is not discarded and it is not replaced by a pointwise smallness claim.
\end{definition}

\begin{definition}[Complete order-one work observable]
\label{def:supp_order_one_work}
Let \(\mathfrak W_{x,p}(q)\) denote the finite-response work observable created
by inserting a full slow-gradient increment \(q(z)\) in the port:
\[
  \mathfrak W_{x,p}(q)
  =
  B(x,p,\cdot)\cdot q
  +(\nabla_\theta\Acal_\theta)[D_p\Theta(x,p)q]\chi_p
  +(\nabla_p B(x,p,\cdot)[q])\cdot p.
\]
Its invariant and centered parts are
\[
  \Lambda_{x,p}(q)=\int\mathfrak W_{x,p}(q)\,d\mu_{\Theta(x,p)},
  \qquad
  \mathfrak W^\circ_{x,p}(q)
  =
  \mathfrak W_{x,p}(q)-\Lambda_{x,p}(q).
\]
For the contact-frozen test, the order-one observable is
\[
  \mathcal G_\varepsilon^\varphi(z)
  =
  V^\circ(x_\varepsilon,p_\varepsilon,z)
  +\mathfrak W_{x_\varepsilon,p_\varepsilon}(q_\varepsilon)(z).
\]
\end{definition}

\begin{proposition}[Order-one average and second cell equation]
\label{prop:supp_second_cell_average}
The invariant part of \(\mathcal G_\varepsilon^\varphi\) is exactly
\[
  \int \mathcal G_\varepsilon^\varphi\,d\mu_{\theta_\varepsilon}
  =
  \Tr(D(x_\varepsilon,p_\varepsilon)X_\varepsilon)
  +H(x_\varepsilon,p_\varepsilon).
\]
The centered part is solved by the deterministic second corrector
\[
  \psi_\varepsilon
  =
  -\Acal_{\theta_\varepsilon}^{-1}
  \Piperp_{\theta_\varepsilon}\mathcal G_\varepsilon^\varphi,
  \qquad
  \Pi_{\theta_\varepsilon}\psi_\varepsilon=0.
\]
Thus the order-one microscopic oscillation is cancelled before the
half-relaxed limit is taken.
\end{proposition}

\begin{proof}
The horizontal component of \(q_\varepsilon\) contributes the finite-response
work term in \(H\).  The Hessian component
\(X_\varepsilon^T\boldsymbol\chi_\varepsilon\) contributes
\[
  \sum_{i,j}X_{\varepsilon,ij}
  \Lambda_{x_\varepsilon,p_\varepsilon}(\chi_i e_j).
\]
The single-action reciprocity of the finite-response port makes the connection
part skew in \(i,j\).  Since \(X_\varepsilon\) is symmetric, the skew part
vanishes after contraction and the remaining symmetric part is the Green-Kubo
tensor.  The residual centered part has zero invariant projection and is in
the range of the reduced observable generator by the Paper 1 resolvent theorem.
\end{proof}

\begin{definition}[Full corrector hierarchy window]
\label{def:supp_full_corrector_hierarchy}
For a contact jet \((x,p,X)\), define the hierarchy by the recursive rule
\[
  \Acal_{\Theta(x,p)}\chi^{(1)}_{a}
  =
  -\Piperp_{\Theta(x,p)}B_a(x,p,\cdot),
  \qquad
  \Pi_{\Theta(x,p)}\chi^{(1)}_a=0,
\]
and, after all order-\(k\) invariant averages have been separated,
\[
  \Acal_{\Theta(x,p)}\chi^{(k+1)}
  =
  -\Piperp_{\Theta(x,p)}
  \mathcal G^{(k)}(x,p,X,\chi^{(1)},\ldots,\chi^{(k)}),
  \qquad
  \Pi_{\Theta(x,p)}\chi^{(k+1)}=0 .
\]
For the HJB limit in the main paper, only \(k=1\) and \(k=2\) are needed.  The
recursive notation records that every higher remainder used for Taylor
bookkeeping would be obtained by the same reduced resolvent on the zero-mean
observable space.
\end{definition}

\begin{proposition}[Hierarchy closure and branch summation]
\label{prop:supp_corrector_hierarchy_closure}
On compact contact-jet windows, every corrector and every differentiated
corrector appearing in \cref{def:supp_full_corrector_hierarchy} is a finite
sum of reduced-resolvent words applied to branchwise finite-response
observables and moving-trace letters.  For each fixed order \(k\le r-8\),
\[
  \|\chi^{(k)}\|_{\#}
  +\sum_{|\alpha|\le r-8-k}
  \|\partial_{x,p}^{\alpha}\chi^{(k)}\|_{w}
  \le C_k(K).
\]
The constants are obtained branchwise and then summed over homogeneity
children by the same growth/strip-tail mechanism as in
\cref{lem:supp_branch_residual_summation}.
\end{proposition}

\begin{proof}
The datum \(\mathcal G^{(k)}\) is built from products of finite-response
read-outs, previous correctors, and horizontal derivatives.  Product bounds in
the anisotropic observable space put its centered part in the domain of the
reduced generator.  Applying the reduced resolvent gives \(\chi^{(k+1)}\).
When a parameter derivative hits the projector, generator, roof, or singular
branch selector, the resolvent identity gives an ordered word with one
differentiated letter.  Interior letters are branchwise smooth; boundary
letters are exactly the trace functionals of
\cref{prop:supp_moving_branch_boundary_ledger}.  The high-strip polynomial
loss is fixed at each finite order, and the admissible strip tail makes the
global child sum finite.  This proves the displayed bounds by induction on
\(k\).
\end{proof}

\section{Residual Identity and Scale Ledger}

\begin{lemma}[Residual identity with normalized gauge]
\label{lem:supp_residual_identity}
Let
\[
  \Phi^\varepsilon
  =
  \varphi
  +\varepsilon\Xi_\varepsilon[\varphi]
  +\varepsilon^2\psi_\varepsilon
\]
and normalize it by subtracting the microscopic gauge value at the contact
point.  When \(\Phi^\varepsilon\) is inserted into the prelimit
viscosity-duality operator and paired with any admissible density
\(\rho\in\Adm_M\), the result is
\[
  \partial_t\varphi(t_\varepsilon,x_\varepsilon)
  +\Hcal(x_\varepsilon,p_\varepsilon,X_\varepsilon)
  +\langle R_\varepsilon[\varphi],\rho\rangle ,
\]
where
\[
  \Hcal(x,p,X)=-\Tr(D(x,p)X)-H(x,p)
\]
and, on compact contact-jet windows,
\[
  |\langle R_\varepsilon[\varphi],\rho\rangle|
  \le
  C_M\bigl(
  \varepsilon
  +|t_\varepsilon-t_0|
  +|x_\varepsilon-x_0|
  +|p_\varepsilon-p_0|
  +|X_\varepsilon-X_0|
  \bigr).
\]
\end{lemma}

\begin{proof}
The scale-\(\varepsilon^{-1}\) terms cancel by
\cref{lem:supp_first_order_cancellation}.  The invariant order-one part is
identified by \cref{prop:supp_second_cell_average}, and its centered part is
cancelled by the second cell equation.  Every remaining term has one of the
following forms:
\begin{center}
\small
\begin{tabular}{p{0.27\linewidth}p{0.38\linewidth}p{0.22\linewidth}}
\hline
Source & Analytic form & Bound \\
\hline
Smooth test expansion &
third-order Taylor and contact displacement terms &
\(C\) times contact displacement \\
First corrector remainder &
\(\varepsilon\partial_t\Xi_\varepsilon\) and horizontal products &
\(C_M\varepsilon\) \\
Second corrector remainder &
\(\varepsilon B\cdot D_x\psi_\varepsilon\) and time derivatives &
\(C_M\varepsilon\) \\
Moving singular boundaries &
trace distributions from changing branch charts &
\(C_M\varepsilon\) after strip-tail summation \\
Parameter continuity &
replacement of the actual contact jet by the limiting jet after cancellation &
\(C\) times contact displacement \\
\hline
\end{tabular}
\end{center}
Product, trace, and differentiated-resolvent bounds are exactly the Paper 1
inputs.  Summing the displayed finite list proves the estimate.
\end{proof}

\begin{proposition}[Four-scale deterministic ledger]
\label{prop:supp_four_scale_ledger}
The perturbed-test calculation is accounted for by
\[
\begin{array}{ll}
\varepsilon^{-2}:&
  \Acal_\theta\varphi=0\quad(\varphi\text{ has no fast variable}),\\[2mm]
\varepsilon^{-1}:&
  \Acal_\theta\chi_p+B\cdot p=0,\\[2mm]
1:&
  V^\circ+\mathfrak W(q)
  \text{ has mean }\Tr(DX)+H
  \text{ and centered part solved by }\psi,\\[2mm]
\varepsilon:&
  \text{Taylor, horizontal, roof, and moving-trace remainders}.
\end{array}
\]
The order-one line uses the actual complete gradient increment
\[
  q=D_x^{\rm fr}\chi_p+X^T\boldsymbol\chi .
\]
Therefore neither a hidden
\(\varepsilon^{-1}\) parameter mismatch nor an omitted
\(D_p\chi\,X\) chain-rule contribution remains.
\end{proposition}

\begin{proof}
The ledger is the combination of
\cref{lem:supp_first_order_cancellation,def:supp_contact_frozen_test,prop:supp_second_cell_average,lem:supp_residual_identity}.
The crucial point is the order of operations: freeze the actual contact jet,
cancel the singular powers, split the order-one term into invariant and
centered parts, and only then pass to the limiting jet.
\end{proof}

\section{Branch Residual Constants}

\begin{definition}[Residual constant ledger]
\label{def:supp_residual_constant_ledger}
Fix an admissible density radius \(M\) and a compact contact-jet window \(K\).
The residual estimate in \cref{lem:supp_residual_identity} uses the following
constants:
\begin{center}
\small
\begin{tabular}{p{0.22\linewidth}p{0.62\linewidth}}
\hline
Constant & Role \\
\hline
\(C_{\rm fr}(K)\) & bounds finite-response derivatives of
\(\Theta,B,V\) on \(K\). \\
\(C_{\rm res}(K)\) & bounds the reduced flow resolvent and the first and second
cell correctors. \\
\(C_{\rm word}(K)\) & bounds differentiated resolvent words generated by
parameter and contact-jet derivatives. \\
\(C_{\rm prod}(M,K)\) & bounds products of admissible densities with
branchwise correctors and horizontal derivatives. \\
\(C_{\rm tr}(M,K)\) & bounds moving-singularity trace letters in the weak
anisotropic norm. \\
\(C_{\rm tail}(M,K,N)\) & bounds the high-strip tail after strip cutoff \(N\),
with \(C_{\rm tail}(M,K,N)\to0\) as \(N\to\infty\). \\
\(C_{\rm ct}(K)\) & bounds the contact displacement modulus for
\((t,x,p,X)\) in \(K\). \\
\hline
\end{tabular}
\end{center}
Choose \(N=N(M,K,\eta)\) so that \(C_{\rm tail}(M,K,N)\le\eta\), and then
choose \(\varepsilon\) small enough that all central-strip, Taylor, and
corrector remainders are below \(\eta\).  The resulting residual modulus is
\[
  \Omega_{M,K}(\varepsilon,\Delta_{\rm ct})
  =
  C_0(M,K)\varepsilon
  +C_{\rm ct}(K)\Delta_{\rm ct}
  +C_{\rm tail}(M,K,N(\varepsilon)),
\]
where
\[
  \Delta_{\rm ct}
  =
  |t_\varepsilon-t_0|+|x_\varepsilon-x_0|
  +|p_\varepsilon-p_0|+|X_\varepsilon-X_0|.
\]
For the contact points selected in the half-relaxed argument,
\(\Omega_{M,K}(\varepsilon,\Delta_{\rm ct})\to0\).
\end{definition}

\begin{lemma}[Branch residual summation]
\label{lem:supp_branch_residual_summation}
Every residual term in the normalized perturbed-test calculation is a finite
sum of homogeneous-branch contributions bounded by the constants in
\cref{def:supp_residual_constant_ledger}.  After summing over branch children,
\[
  |\langle R_\varepsilon[\varphi],\rho\rangle|
  \le
  \Omega_{M,K}(\varepsilon,\Delta_{\rm ct})
  (1+\|\varphi\|_{C^4})
\]
for all \(\rho\in\Adm_M\).
\end{lemma}

\begin{proof}
On a central homogeneous branch, the Taylor remainder and the finite-response
coefficient remainder are ordinary \(C^4\) estimates with constants
\(C_{\rm fr}\) and \(C_{\rm ct}\).  The first and second corrector remainders
are products of admissible densities, correctors, and horizontal derivatives;
they are bounded by \(C_{\rm prod}C_{\rm res}\varepsilon\).  Differentiating
the parameter-dependent projector or resolvent gives finitely many ordered
resolvent words, bounded by \(C_{\rm word}\).

When the branch cut moves, the residual is a weak trace functional supported on
a singular hypersurface.  The trace bound \(C_{\rm tr}\) controls its central
strip contribution.  In high strips, the crossing denominator grows only as a
fixed polynomial in the strip index, while admissible densities carry the
strip-tail weight supplied by the response theorem.  Hence the total high-strip
contribution is \(C_{\rm tail}(M,K,N)\).  Summing first over central branches
and then over the high-strip tail gives the displayed bound.
\end{proof}

\begin{definition}[Branch residual family decomposition]
\label{def:supp_branch_residual_family}
For a fixed admissible-density radius \(M\), compact contact-jet window \(K\),
and strip cutoff \(N\), decompose the normalized residual as
\[
  R_\varepsilon
  =
  R_\varepsilon^{\rm cent}
  +R_\varepsilon^{\rm word}
  +R_\varepsilon^{\rm tr}
  +R_\varepsilon^{\rm end}
  +R_\varepsilon^{\rm tail}.
\]
The five terms are defined by their source after freezing the contact jet:
\begin{center}
\small
\begin{tabular}{p{0.2\linewidth}p{0.66\linewidth}}
\hline
Term & Source \\
\hline
\(R_\varepsilon^{\rm cent}\) &
central-branch Taylor, roof, and smooth finite-response coefficient
remainders. \\
\(R_\varepsilon^{\rm word}\) &
differentiated projector, generator, reduced-resolvent, and corrector words
with all letters supported in the central strip family. \\
\(R_\varepsilon^{\rm tr}\) &
moving grazing, homogeneity-boundary, and branch-change trace letters in
central strips. \\
\(R_\varepsilon^{\rm end}\) &
endpoint and terminal-gauge terms, including specular endpoint variations and
cash-additive terminal shifts. \\
\(R_\varepsilon^{\rm tail}\) &
all central estimates repeated in strips above \(N\), plus their terminal
children. \\
\hline
\end{tabular}
\end{center}
This decomposition is made after the scale-\(\varepsilon^{-1}\) cancellation
and after the order-one invariant average has been separated.  Hence no term
in the table changes the effective coefficients; the table only audits the
vanishing residual.
\end{definition}

\begin{lemma}[Central branch residual check]
\label{lem:supp_central_branch_residual_check}
For each central homogeneous branch \(b\) below the cutoff \(N\),
\[
  \bigl|\langle
  R_{\varepsilon,b}^{\rm cent}
  +R_{\varepsilon,b}^{\rm word}
  +R_{\varepsilon,b}^{\rm end},\rho\rangle\bigr|
  \le
  C_{M,K,N}
  (1+\|\varphi\|_{C^5})
  (\varepsilon+\Delta_{\rm ct})
\]
uniformly for \(\rho\in\Adm_M\).  The constant depends on \(N\), but not on
\(\varepsilon\), the branch label \(b\), or the selected upper/lower contact.
\end{lemma}

\begin{proof}
On a central branch the collision map, free-flight roof, finite-response
selector, and branch Jacobian are \(C^{r-6}\) functions of the contact jet.
The Taylor remainder of the branch principal action is therefore bounded by
\(C_{K,N}\|\varphi\|_{C^5}(\varepsilon+\Delta_{\rm ct})\).  If a derivative
hits the invariant projector, reduced generator, or corrector, the resolvent
identity produces a finite ordered word.  Every letter is central and has
bounded strong or weak norm on \(K\), so the product bound gives the same
factor.  Endpoint terms are handled by
\cref{lem:supp_specular_endpoint_variation,prop:supp_terminal_stability}: the
physical reflection contributes no extra boundary work, while a terminal
constant changes only the cash-additive level and not the residual
derivatives.  There are finitely many central branch families below \(N\),
which gives the displayed uniform constant.
\end{proof}

\begin{lemma}[Trace branch residual check]
\label{lem:supp_trace_branch_residual_check}
The central-strip moving-boundary part satisfies
\[
  |\langle R_\varepsilon^{\rm tr},\rho\rangle|
  \le
  C_{\rm tr}(M,K,N)
  (1+\|\varphi\|_{C^5})
  (\varepsilon+\Delta_{\rm ct})
\]
for every admissible density \(\rho\in\Adm_M\).  The trace constant is the same
for upper and lower half-relaxed contacts.
\end{lemma}

\begin{proof}
Moving singular hypersurfaces enter only when differentiating a branch
selector or a strip cutoff.  In local normal coordinates the derivative is a
coarea trace with a normal crossing denominator.  The finite-response port
keeps the crossing speed and the contact jet in a compact window, so the
central denominator loss is bounded by \(C_{\rm tr}(M,K,N)\).  The trace is a
weak anisotropic functional, exactly the class controlled by the imported
moving-response theorem.  Upper and lower contacts use the same branch
selectors and the same normalized gauge; only the sign of the final viscosity
inequality changes, not the trace norm.
\end{proof}

\begin{proposition}[High-strip residual cutoff]
\label{prop:supp_high_strip_residual_cutoff}
For every \(\eta>0\) there is a cutoff \(N=N(M,K,\eta)\) such that
\[
  |\langle R_\varepsilon^{\rm tail},\rho\rangle|
  \le
  \eta(1+\|\varphi\|_{C^5})
\]
for all sufficiently small \(\varepsilon\), uniformly over all contacts in
\(K\) and all \(\rho\in\Adm_M\).
\end{proposition}

\begin{proof}
In strip \(n\), differentiating the branch action can lose only the fixed
polynomial power already listed in the Paper 1 response budget: impact time,
roof, branch Jacobian, trace multiplicity, and the finite number of corrector
word letters.  The admissible density class carries a strip-tail weight with a
larger exponent.  Summing over terminal descendants first gives a geometric
standard-pair factor, while the remaining polynomial strip sum tends to zero
as \(N\to\infty\).  The choice of \(N\) is made before sending
\(\varepsilon\to0\), so the cutoff is uniform for the half-relaxed passage.
\end{proof}

\begin{proposition}[Full branch residual check]
\label{prop:supp_full_branch_residual_check}
With the decomposition of \cref{def:supp_branch_residual_family}, the residual
modulus in \cref{lem:supp_branch_residual_summation} may be chosen as
\[
  \Omega_{M,K}(\varepsilon,\Delta_{\rm ct})
  =
  C_{M,K,N}(1+\|\varphi\|_{C^5})(\varepsilon+\Delta_{\rm ct})
  +\eta(1+\|\varphi\|_{C^5}),
\]
where \(N=N(M,K,\eta)\).  Consequently, after first choosing \(N\) and then
letting \(\varepsilon\to0\), the branch residual vanishes uniformly in the
subsolution and supersolution passages.
\end{proposition}

\begin{proof}
Add the five terms in \cref{def:supp_branch_residual_family}.  Central smooth,
word, and endpoint terms are controlled by
\cref{lem:supp_central_branch_residual_check}.  Central moving-boundary traces
are controlled by \cref{lem:supp_trace_branch_residual_check}.  The
high-strip remainder is controlled by
\cref{prop:supp_high_strip_residual_cutoff}.  These bounds are uniform over
the paired upper/lower contacts used in \cref{lem:supp_contact_selection};
therefore the order of limits gives the displayed modulus and the vanishing
residual required by \cref{prop:supp_subsuper_passage}.
\end{proof}

\begin{proposition}[Uniform sub/sup residual modulus]
\label{prop:supp_uniform_subsuper_modulus}
The subsolution and supersolution inequalities in
\cref{prop:supp_subsuper_passage} hold with the same residual modulus
\(\Omega_{M,K}\).  The sign of the viscosity inequality changes between the
two halves, but the constants and the branch summations are identical.
\end{proposition}

\begin{proof}
Both upper and lower contacts use the same normalized gauge, the same
contact-frozen first corrector, and the same second cell equation.  The
branchwise prelimit operator is the same deterministic action operator in both
cases.  Therefore the residual sources listed in
\cref{def:supp_residual_constant_ledger} are unchanged.  Only the inequality
direction in the viscosity-duality test changes, so the same
\(\Omega_{M,K}\) controls both limits.
\end{proof}

\section{Smooth Approximation in the Weak Topology}

\begin{lemma}[Admissible density smoothing]
\label{lem:supp_admissible_density_smoothing}
Fix an admissible-density ball \(\Adm_M\), a compact contact-jet window \(K\),
and \(\eta>0\).  For every \(\rho\in\Adm_M\) there is a density
\(\rho_{\eta}\) that is \(C^{r-6}\) on each homogeneous rectangle, vanishes in
an \(\eta\)-neighborhood of the canonical singular strata and above a finite
strip cutoff, has total mass one, and satisfies
\[
  \|\rho_\eta\|_{\Adm_{C M}}\le C M,\qquad
  |\langle F,\rho-\rho_\eta\rangle|
  \le \eta\,\|F\|_{\rm test}
\]
for every branchwise test observable \(F\) appearing in the HJB residual
ledger.  The constant \(C\) is independent of \(\rho\) and \(\eta\).
\end{lemma}

\begin{proof}
First remove high homogeneity strips, choosing the cutoff so that the
admissible strip tail contributes at most \(\eta/3\).  Next remove an
\(\eta\)-thin neighborhood of the finite singular-stratum family below the
cutoff.  Its contribution is controlled by the weak trace bound used in
\cref{lem:supp_trace_branch_residual_check}.  On the remaining finite family
of homogeneous rectangles, mollify in collision coordinates along stable and
unstable chart directions and use a \(C^{r-6}\) partition of unity.  The
anisotropic density norm is stable under this chartwise smoothing, up to a
constant depending only on the finite atlas.  Renormalization changes pairings
by \(O(\eta)\), because the removed mass is \(O(\eta)\) uniformly on the
admissible ball.  This proves the displayed bounds.
\end{proof}

\begin{proposition}[Perturbed-test smoothing]
\label{prop:supp_perturbed_test_smoothing}
Let \(\Phi^\varepsilon\) be the normalized contact-frozen perturbed test used
in \cref{lem:supp_residual_identity}.  For every compact contact-jet window
\(K\), admissible-density radius \(M\), and \(\eta>0\), there is a branchwise
\(C^{r-8}\) perturbed test \(\Phi^\varepsilon_\eta\), supported away from a
finite singular-neighborhood and using the same frozen macroscopic jet, such
that
\[
  \sup_{\rho\in\Adm_M}
  \bigl|
  \langle
  \Lcal_\varepsilon \Phi^\varepsilon
  -\Lcal_\varepsilon \Phi^\varepsilon_\eta,\rho
  \rangle
  \bigr|
  \le
  \eta(1+\|\varphi\|_{C^5})
\]
for all sufficiently small \(\varepsilon\).  The subsolution and supersolution
inequalities are unchanged after sending \(\eta\downarrow0\).
\end{proposition}

\begin{proof}
The correctors in \(\Phi^\varepsilon\) are reduced-resolvent images of
branchwise finite-response observables and moving-trace letters.  Approximate
each observable by the chartwise smoothing procedure in
\cref{lem:supp_admissible_density_smoothing}; apply the reduced resolvent to
the smoothed zero-mean observable; and subtract the invariant mean to preserve
the normalized gauge.  Differentiated resolvent words commute with this
approximation up to the same weak error, because their letters are bounded by
the response package and because the singular-neighborhood contribution is
already a trace term.  The prelimit operator \(\Lcal_\varepsilon\) contains
only finitely many branch derivatives on the chosen order window, so the weak
pairing error is bounded by the displayed \(\eta\)-term.  The viscosity
inequalities are first proved for the smoothed tests and then recovered by
letting \(\eta\to0\).
\end{proof}

\begin{corollary}[Smooth approximation closure]
\label{cor:supp_smooth_approximation_closure}
The detailed HJB closure theorem is unchanged if all admissible densities and
all microscopic correctors used in the perturbed-test argument are first
approximated by branchwise smooth objects supported away from singular strata,
with the approximation error sent to zero before the half-relaxed limit.
\end{corollary}

\begin{proof}
Combine \cref{lem:supp_admissible_density_smoothing,prop:supp_perturbed_test_smoothing}.
The order of limits is: choose a compact contact window, choose the smoothing
and strip cutoff so the weak approximation error is below \(\eta\), send
\(\varepsilon\to0\) in the smoothed prelimit inequality, and finally send
\(\eta\to0\).  The branch residual modulus of
\cref{prop:supp_full_branch_residual_check} is uniform in this procedure.
\end{proof}

\section{Half-Relaxed Limits}

\begin{lemma}[Contact selection for normalized gauges]
\label{lem:supp_contact_selection}
Let \(\varphi\in C^3\) touch the upper half-relaxed paired limit strictly at
\((t_0,x_0)\).  Then there are prelimit contact points
\((t_\varepsilon,x_\varepsilon)\to(t_0,x_0)\) and admissible densities
\(\rho_\varepsilon\) such that the normalized contact-frozen perturbed test is
an upper viscosity-duality test for \(u^\varepsilon\).  The lower-contact
statement holds with inequalities reversed.
\end{lemma}

\begin{proof}
Add a small strict macroscopic bump and work on a compact cylinder.  The paired
half-relaxed definition gives approximate maximizing triples.  The microscopic
gauge is normalized by subtracting its value at the chosen contact point; this
does not alter generator or horizontal derivatives, but it restores the exact
contact equality in the paired test.  Homogeneous branch approximations avoid
the zero-measure iterated singular strata, and high-strip tails are summable by
the response-theory trace estimates.  A deterministic bump turns approximate
contacts into exact dual contacts.  The lower-contact case is identical.
\end{proof}

\begin{proposition}[Subsolution and supersolution passage]
\label{prop:supp_subsuper_passage}
The upper half-relaxed paired envelope is a viscosity subsolution of
\[
  \partial_t u+\Hcal(x,Du,D^2u)=0,
\]
and the lower half-relaxed paired envelope is a viscosity supersolution.  The
same first and second cell equations are used in both halves; only the
prelimit viscosity inequality changes sign.
\end{proposition}

\begin{proof}
At an upper contact, use \cref{lem:supp_contact_selection} and insert the
normalized perturbed test in the prelimit subsolution inequality.  By
\cref{lem:supp_residual_identity},
\[
  \partial_t\varphi(t_\varepsilon,x_\varepsilon)
  +\Hcal(x_\varepsilon,D\varphi,D^2\varphi)
  \le
  |\langle R_\varepsilon[\varphi],\rho_\varepsilon\rangle|.
\]
The residual bound tends to zero on compact contact-jet windows, giving the
subsolution inequality.  The lower contact gives the reverse inequality for
the lower envelope by the same argument.
\end{proof}

\begin{lemma}[Density-to-state upgrade]
\label{lem:supp_density_to_state}
Suppose \(F_\varepsilon(z)\) is uniformly bounded on homogeneous rectangles,
its pairings against every fixed admissible density ball tend to zero, and its
oscillation on nested homogeneous subrectangles is bounded by
\[
  \operatorname{osc}_{R'}F_\varepsilon
  \le \omega_R(\operatorname{diam}R')+r_\varepsilon,
  \qquad r_\varepsilon\to0.
\]
Then \(F_\varepsilon(z)\to0\) for Liouville-a.e. microscopic state after
discarding the canonical iterated singular set.
\end{lemma}

\begin{proof}
Disintegrate Liouville measure into a countable algebra of homogeneous
rectangles.  A stable-smoothed density equal to the normalized indicator of a
fixed rectangle, up to an arbitrarily small boundary layer, has a finite
anisotropic norm.  Hence the pairing hypothesis applies to that rectangle.
If a positive-measure set had \(\limsup |F_\varepsilon|>\eta\), the leafwise
differentiation basis and the oscillation modulus would give a smaller fixed
rectangle with an average of one sign bounded away from zero, contradicting
the pairing limit.  Countability gives the almost-everywhere statement.
\end{proof}

\begin{theorem}[Detailed HJB closure]
\label{thm:supp_detailed_hjb_closure}
The prelimit deterministic viscosity-duality solutions converge, after
anisotropic microscopic pairing and then almost everywhere in the homogeneous
state variable, to the unique viscosity solution of the effective HJB equation.
The proof uses only the finite-response port, the imported Paper 1 response
theorem, the corrector hierarchy above, and macroscopic viscosity comparison.
\end{theorem}

\begin{proof}
The upper and lower paired envelopes inherit terminal data from the prelimit
terminal condition.  By \cref{prop:supp_subsuper_passage} they are respectively
sub- and supersolutions of the effective HJB.  The comparison modulus recorded
in the main article gives comparison, so the envelopes coincide with the unique
viscosity solution.  Uniformity of \cref{lem:supp_residual_identity} over
bounded admissible density balls gives paired convergence.  Applying
\cref{lem:supp_density_to_state} on a countable dense family of macroscopic
points and then using deterministic macroscopic equicontinuity gives the
Liouville-a.e. state statement.
\end{proof}

\section{Coefficient and Nonconvexity Checks}

\begin{proposition}[Coefficient differentiability ledger]
\label{prop:supp_coefficient_differentiability}
On every compact contact-jet window, the maps
\[
  (x,p)\mapsto D(x,p),\qquad (x,p)\mapsto H(x,p)
\]
have the derivatives used in the comparison and residual estimates.  Each
derivative is a finite sum of invariant projector derivatives, reduced
resolvent derivatives, primitive read-out derivatives, and moving-boundary
trace insertions.
\end{proposition}

\begin{proof}
Differentiate the corrector identity
\[
  \chi_i=-\Acal_{\Theta(x,p)}^{-1}B_i(x,p,\cdot)
\]
in the moving Banach bundle.  The derivative of the inverse generator is the
usual resolvent word with one differentiated generator letter; higher
derivatives give ordered resolvent words.  Paper 1 bounds those words,
including trace letters caused by moving singular boundaries.  Since the
finite-response read-outs have bounded branchwise derivatives on the compact
window, differentiating the invariant pairings defining \(D\) and \(H\) gives
the stated finite expansion.
\end{proof}

\begin{definition}[Correlation Hilbert bundle]
\label{def:supp_correlation_hilbert_bundle}
For each compact contact window and each \((x,p)\), define a bilinear
correlation form on centered finite-response observables by
\[
  \mathfrak c_{x,p}(f,g)
  =
  \frac12\int_0^\infty
  \int \bigl(f(z)g(\Phi_{\Theta(x,p)}^t z)
      +g(z)f(\Phi_{\Theta(x,p)}^t z)\bigr)
  \,d\mu_{\Theta(x,p)}(z)\,dt .
\]
The integral is interpreted through the reduced observable resolvent.  Its null
space is factored out, and the completion is denoted
\(\mathfrak H_{x,p}\).  The drift factor is the map
\[
  \Sigma(x,p)^*e_i=[B_i(x,p,\cdot)]_{\mathfrak H_{x,p}},
  \qquad
  D(x,p)=\Sigma(x,p)\Sigma(x,p)^* .
\]
\end{definition}

\begin{lemma}[Green-Kubo positivity and coboundary kernel]
\label{lem:supp_gk_coboundary_kernel}
The tensor \(D(x,p)\) is symmetric positive semidefinite.  Moreover,
\[
  a^TD(x,p)a=0
\]
if and only if the centered observable \(a\cdot B(x,p,\cdot)\) is a
deterministic generator coboundary,
\[
  a\cdot B(x,p,\cdot)=-\Acal_{\Theta(x,p)}h
\]
for some zero-mean observable \(h\) in the reduced generator domain.
\end{lemma}

\begin{proof}
Symmetry is built into \(\mathfrak c_{x,p}\).  Positivity follows from the
finite-time variance identity
\[
  a^TD(x,p)a
  =
  \lim_{T\to\infty}\frac1{2T}
  \int\left(\int_0^T a\cdot B(x,p,\Phi^t z)\,dt\right)^2d\mu(z),
\]
where convergence is justified by deterministic correlation summability.  If
\(a\cdot B=-\Acal h\), the time integral telescopes to
\(h(\Phi^Tz)-h(z)\), so the normalized variance vanishes.  Conversely, if the
variance vanishes, the reduced resolvent \(h=-\Acal^{-1}(a\cdot B)\) has zero
correlation energy; the zero-mean spectral gap identifies the observable with
the generator range.  This is an anisotropic observable-domain statement, not
an \(L^2\) mixing assumption.
\end{proof}

\begin{proposition}[Correlation-factor transport]
\label{prop:supp_correlation_factor_transport}
The fibers \(\mathfrak H_{x,p}\) in
\cref{def:supp_correlation_hilbert_bundle} admit deterministic metric
transports \(J^p_{y\to x}:\mathfrak H_{y,p}\to\mathfrak H_{x,p}\) on compact
coordinate windows such that
\[
  \|\,\Sigma(x,p)-J^p_{y\to x}\Sigma(y,p)\,\|_{\mathrm{HS}}
  \le C_K|x-y|.
\]
This supplies the Crandall-Ishii diffusion factor without taking a matrix
square root of a possibly degenerate covariance matrix.
\end{proposition}

\begin{proof}
Differentiate the correlation form in the moving Banach bundle.  The derivative
is a finite sum of invariant-projector, reduced-resolvent, primitive
observable, and moving-boundary trace terms, all bounded by the response
package.  On the quotient by the null space, Riesz representation gives the
unique finite-rank correction that makes the covariant derivative metric:
\[
  \partial_{x_a}\mathfrak c_{x,p}(f,g)
  =
  \mathfrak c_{x,p}(\nabla_a f,g)
  +\mathfrak c_{x,p}(f,\nabla_a g).
\]
Solving the associated linear ODE along the segment from \(y\) to \(x\) gives
the isometric transport \(J^p_{y\to x}\).  Since \(B_i\) is a
finite-response observable with bounded moving-bundle derivative, its
transported class is Lipschitz in the correlation Hilbert norm.  Summing over
the finitely many slow coordinates gives the displayed Hilbert-Schmidt bound.
\end{proof}

\begin{proposition}[Green-Kubo comparison trace estimate]
\label{prop:supp_gk_comparison_trace}
If \(p=\alpha(x-y)\) and \(X,Y\) satisfy the Crandall-Ishii block inequality,
then on compact windows
\[
  -\Tr(D(x,p)X-D(y,p)Y)
  \le C\alpha |x-y|^2+o_{\rm fr}(1),
\]
where \(o_{\rm fr}(1)\) is the finite-rank correlation-bundle tail removed
after the Hilbert projection cutoff.
\end{proposition}

\begin{proof}
Transport \(\Sigma(y,p)\) into \(\mathfrak H_{x,p}\) by
\cref{prop:supp_correlation_factor_transport}.  Apply the Ishii matrix
inequality to the paired vectors
\((\Sigma(x,p)^*\xi,J^p_{y\to x}\Sigma(y,p)^*\xi)\) in a finite-dimensional
subspace of \(\mathfrak H_{x,p}\).  The block inequality gives
\[
  -\Tr(D(x,p)X-D(y,p)Y)
  \le
  3\alpha\|\Sigma(x,p)-J^p_{y\to x}\Sigma(y,p)\|_{\mathrm{HS}}^2
  +o_{\rm fr}(1).
\]
The factor-transport bound gives the stated estimate, and the projection tail
vanishes by monotone convergence of the positive correlation form and the
uniform reduced-resolvent bound.
\end{proof}

\begin{proposition}[Comparison modulus ledger]
\label{prop:supp_comparison_modulus_ledger}
The derived operator
\[
  \Hcal(x,p,X)=-\Tr(D(x,p)X)-H(x,p)
\]
is continuous, degenerate elliptic, and satisfies the Crandall-Ishii modulus on
compact contact-jet windows.
\end{proposition}

\begin{proof}
Degenerate ellipticity follows from positivity of the Green-Kubo tensor.  The
differentiability ledger gives local continuity and Lipschitz control in
\((x,p)\).  The Green-Kubo representation gives a Hilbert-factorization
\(D=\Sigma\Sigma^*\).  Applying the Ishii block inequality to
\((\Sigma(x,p)^*\xi,\Sigma(y,p)^*\xi)\) and using finite-rank approximation in
the correlation Hilbert bundle yields the standard trace estimate.  Combining
that estimate with the local modulus of \(H\) gives the comparison condition.
\end{proof}

\begin{lemma}[Synchronized density for doubled contacts]
\label{lem:supp_synchronized_density_doubling}
In the doubled-variable comparison argument, the same microscopic admissible
density can be used in the upper and lower viscosity inequalities after
transport through the two Banach-bundle charts.  If the two contact fibers are
\(\Theta(x,p_U)\) and \(\Theta(y,p_V)\), then the transported density satisfies
\[
  \|\rho_U-\mathfrak T_{\Theta(x,p_U)\leftarrow\Theta(y,p_V)}\rho_V\|_{w}
  \le
  C_M\bigl(|x-y|+\|p_U-p_V\|_{w}\bigr)+o_{\rm strip}(1),
\]
where \(o_{\rm strip}(1)\) is sent to zero by the common homogeneity cutoff.
\end{lemma}

\begin{proof}
The normal-coordinate transport between nearby billiard fibers preserves
positivity and total mass and has uniformly bounded anisotropic norm on compact
contact windows.  Interior generator, drift, and potential terms differ by the
finite-response Lipschitz modulus.  Physical collision faces have cancelling
incoming and outgoing flux by \cref{lem:supp_specular_endpoint_variation}.
Grazing and branch-change discrepancies are moving-boundary trace terms, hence
are bounded by \cref{prop:supp_moving_branch_boundary_ledger}.  This gives the
displayed modulus and permits a single synchronized density in the doubled
pairing.
\end{proof}

\begin{proposition}[Doubled-variable comparison refinement]
\label{prop:supp_doubled_variable_comparison}
Let \(u\) be a bounded upper semicontinuous subsolution and \(v\) a bounded
lower semicontinuous supersolution of the derived terminal-value HJB, with
ordered terminal data.  On compact cotangent windows the standard
Crandall-Ishii comparison proof closes with the coefficient modulus of
\cref{prop:supp_comparison_modulus_ledger} and the synchronized-density
estimate of \cref{lem:supp_synchronized_density_doubling}.  In particular no
additional microscopic order assumption is needed.
\end{proposition}

\begin{proof}
Maximize
\[
  u(t,x)-v(s,y)-\frac{|x-y|^2}{2\alpha}
  -\frac{|t-s|^2}{2\alpha}-\delta(|x|^2+|y|^2)
\]
on the compact doubled torus.  The terminal ordering puts the maximizing point
in the parabolic interior for small parameters.  Ishii's lemma supplies common
gradient \(p_\alpha=(x-y)/\alpha\) and matrices \(X,Y\) satisfying the usual
block inequality.  The diffusion term is estimated by applying that block
inequality to a finite-rank approximation of the Green-Kubo correlation
factor.  The error is then removed in the correlation Hilbert bundle.  The
Hamiltonian part is controlled by the finite-response modulus in \(x\) and
\(p\).

For the prelimit paired comparison used in the half-relaxed proof, the upper
and lower microscopic tests are first transported to one reference fiber.
\Cref{lem:supp_synchronized_density_doubling} bounds the density-transport
error, while \cref{prop:supp_moving_branch_boundary_ledger} removes the
singular trace tail.  Subtracting the two viscosity inequalities leaves
\[
  0\le o_\alpha(1)+o_\delta(1)
\]
at a hypothetical positive maximum, a contradiction after sending
\(\alpha\downarrow0\) and then \(\delta\downarrow0\).  Perron stability and
terminal-data Lipschitz dependence follow by the same comparison with constant
barriers, using the zero-energy gauge \(H(x,0)=0\).
\end{proof}

\begin{proposition}[Terminal-value stability and zero-energy gauge]
\label{prop:supp_terminal_stability}
The effective terminal-value HJB has a unique bounded uniformly continuous
viscosity solution and the solution operator \(S_t\phi=u(0,\cdot)\) satisfies
\[
  \|S_t\phi-S_t\psi\|_\infty\le \|\phi-\psi\|_\infty,\qquad
  S_t(\phi+c)=S_t\phi+c .
\]
The same statements hold for the paired prelimit problem on each fixed
\(\varepsilon>0\) in the admissible-density order.
\end{proposition}

\begin{proof}
Comparison from \cref{prop:supp_doubled_variable_comparison} gives uniqueness.
For stability, compare the solution with terminal datum \(\phi\) to the
solutions with terminal data \(\psi+\|\phi-\psi\|_\infty\) and
\(\psi-\|\phi-\psi\|_\infty\).  Constants solve the homogeneous equation
because the zero-energy normalization gives \(H(x,0)=0\) and the diffusion term
vanishes on constants.  The prelimit paired statement is identical: the
branchwise Hamilton principal function is shifted by the terminal constant,
the microscopic generator kills constants, and the paired order is separated
by admissible densities.
\end{proof}

\begin{proposition}[Concrete parameter window for the two-mode port]
\label{prop:supp_concrete_parameter_window}
For the triangular Lorentz cell and the two disjoint deformation bumps used in
the main article, there is a nonempty open parameter window in which the
finite-horizon geometry remains unchanged and the nonconvexity dominance
condition holds.  Let \(\theta_*<\theta_{\rm geom}\) be small enough to
preserve curvature, separation, and corridor blocking for all normal
displacements with \(C^{r+3}\)-norm at most \(\theta_*\).  Choose the action
coefficient \(\lambda(x)\), quartic stabilizer \(\kappa(x)>0\), and response
amplitudes so that
\[
  2\lambda(x)\langle\eta\rangle_{\theta_0}
  >
  C_{\rm resp}(x)+2\delta_x,
  \qquad
  \|\Theta(x,p)-\theta_0\|_{C^{r+3}}\le \theta_*/2
\]
on the working cotangent window.  Then every sufficiently small perturbation
of the port coefficients inside this window still satisfies the strict
dominance inequality with margin at least \(\delta_x\).
\end{proposition}

\begin{proof}
The finite-horizon triangular Lorentz inequalities are open in the
\(C^{r+3}\) topology of the scatterer boundary, so a geometric radius
\(\theta_{\rm geom}>0\) exists.  The normalized deformation bumps and the
observable bump have disjoint supports, hence small amplitudes keep the table
in the same finite-horizon class and avoid branch interference among the two
response modes.

The response remainder \(C_{\rm resp}(x)\) is continuous in the primitive
read-out norms and in the Paper 1 response constants over the compact working
window.  The observable \(\eta\) is nonnegative and nonzero on a regular
collision arc, so \(\langle\eta\rangle_{\theta_0}>0\).  Increasing
\(\lambda(x)\) within the finite-response energy budget, or first reducing the
response amplitudes and then fixing \(\lambda\), gives the strict margin
\(2\delta_x\).  Continuity leaves margin \(\delta_x\) under nearby parameter
choices.  The positive quartic term keeps the example inside the compact
cotangent window and does not remove the negative second directional
coefficient at the origin.
\end{proof}

\begin{corollary}[Concrete parameter-choice recipe]
\label{cor:supp_concrete_parameter_recipe}
The nonconvex port can be chosen by the following finite sequence.
\begin{enumerate}[label=\textup{(\roman*)},wide,labelindent=0pt]
  \item Fix the triangular Lorentz cell and choose the two disjoint normalized
  boundary bumps and the nonnegative observable bump \(\eta\).
  \item Choose a geometric radius \(\theta_*\) inside the finite-horizon
  window, and choose response amplitudes so that
  \(\|\Theta(x,p)-\theta_0\|_{C^{r+3}}\le\theta_*/2\).
  \item Compute the response-budget bound \(C_{\rm resp}(x)\) from the imported
  Paper 1 constants on the compact cotangent window.
  \item Choose \(\lambda(x)\) so that
  \(2\lambda(x)\langle\eta\rangle_{\theta_0}>C_{\rm resp}(x)+2\delta_x\).
  \item Choose a positive quartic stabilizer \(\kappa(x)\) large enough to keep
  the action window compact without changing the negative second directional
  coefficient at \(p=0\).
\end{enumerate}
Every inequality in this list is open; hence nearby smooth parameter choices
give the same nonconvexity certificate.
\end{corollary}

\begin{proof}
The first two steps put the table family in the geometric window.  The third
step is a finite evaluation of already imported response constants.  Since
\(\eta\ge0\) and is not identically zero on a regular collision arc, its
invariant average is positive, so the fourth step is possible.  The quartic
term has zero second derivative at the origin in the chosen direction after
normalization and therefore affects confinement but not the local nonconvex
coefficient.  Openness follows from continuity of the finite-response read-outs
and of the response-budget constants over the compact window.
\end{proof}

\begin{definition}[Dimensionless parameter audit table]
\label{def:supp_dimensionless_parameter_audit}
For a compact cotangent window normalize the concrete-port quantities by
\[
  \mathcal R
  =
  \frac{2\lambda\langle\eta\rangle_{\theta_0}}{C_{\rm resp}+2\delta_x},
  \qquad
  \mathcal A
  =
  \frac{\|\Theta-\theta_0\|_{C^{r+3}}}{\theta_*},
  \qquad
  \mathcal S
  =
  \frac{\kappa}{1+\lambda+\|\Theta-\theta_0\|_{C^{r+3}}}.
\]
The following dimensionless choices are sufficient for the port used in the
main paper:
\begin{center}
\small
\begin{tabular}{p{0.28\linewidth}p{0.22\linewidth}p{0.38\linewidth}}
\hline
Quantity & Audit value & Consequence \\
\hline
geometric amplitude ratio \(\mathcal A\) & \(\le 1/2\) &
the table stays in the finite-horizon billiard window. \\
nonconvex dominance ratio \(\mathcal R\) & \(\ge 2\) &
the negative direct action term dominates all response-corrector terms. \\
quartic stabilizer ratio \(\mathcal S\) & \(\ge 4\) &
the action remains confined on the selected cotangent window. \\
smoothing radius \(\delta_{\rm sm}\) & \(\le 1/16\) of the bump separation &
the two boundary bumps and the observable bump remain disjoint. \\
strip-tail tolerance & \(\le 10^{-3}\delta_x\) &
the high-strip contribution cannot close the nonconvexity gap. \\
\hline
\end{tabular}
\end{center}
This is an audit table for normalized inequalities, not a numerical simulation
of a billiard orbit.
\end{definition}

\begin{proposition}[Parameter audit verification]
\label{prop:supp_parameter_audit_verification}
If the dimensionless inequalities in
\cref{def:supp_dimensionless_parameter_audit} hold on the compact cotangent
window, then the concrete two-mode port satisfies the hypotheses of
\cref{prop:supp_concrete_parameter_window,cor:supp_concrete_parameter_recipe}
after the smoothing and strip cutoffs of
\cref{cor:supp_smooth_approximation_closure}.
\end{proposition}

\begin{proof}
The amplitude bound \(\mathcal A\le1/2\) leaves half of the finite-horizon
geometric margin unused, so the smoothed port cannot create a new grazing
family.  The dominance bound \(\mathcal R\ge2\) gives
\[
  2\lambda\langle\eta\rangle_{\theta_0}
  \ge 2(C_{\rm resp}+2\delta_x),
\]
so even after losing \(C_{\rm resp}\) to response-corrector terms and
\(\delta_x\) to smoothing, a gap of at least \(\delta_x\) remains.  The
stabilizer bound keeps the finite-response action coercive on the chosen
cotangent window and does not change the second derivative at the origin in
the nonconvex direction.  The bump-separation condition preserves the sign and
support assumptions used in the direct action computation.  Finally, the
strip-tail tolerance is smaller than the remaining gap, so
\cref{prop:supp_high_strip_residual_cutoff} cannot erase the certificate.
Thus all inequalities in the concrete parameter recipe remain valid after the
regularization steps.
\end{proof}

\begin{definition}[Normalized numerical certificate]
\label{def:supp_normalized_numerical_certificate}
The following dimensionless values give an inspectable certificate for the
concrete-port inequalities.  Work in units in which
\[
  C_{\rm resp}=1,\qquad
  \langle\eta\rangle_{\theta_0}=1,\qquad
  \theta_*=1,
  \qquad
  b_{\rm sep}=1,
\]
where \(b_{\rm sep}\) is the minimum separation among the two boundary bumps
and the observable bump.  Choose
\[
  \delta_x=\frac1{10},\quad
  \lambda=\frac32,\quad
  \|\Theta-\theta_0\|_{C^{r+3}}=\frac25,\quad
  \kappa=16,\quad
  \delta_{\rm sm}=\frac1{32},\quad
  \tau_{\rm strip}=5\cdot 10^{-5}.
\]
Then
\[
  \mathcal A=\frac25,\qquad
  \mathcal R=\frac{3}{6/5}=\frac52,\qquad
  \mathcal S=\frac{16}{1+3/2+2/5}=\frac{160}{29}>5,
\]
and
\[
  \delta_{\rm sm}<\frac{b_{\rm sep}}{16},
  \qquad
  \tau_{\rm strip}<10^{-3}\delta_x .
\]
This is a numerical illustration of the normalized inequalities in
\cref{def:supp_dimensionless_parameter_audit}; it is not an orbit simulation
and it does not estimate the Paper 1 response constants from trajectory data.
\end{definition}

\begin{proposition}[Numerical certificate margin]
\label{prop:supp_numerical_certificate_margin}
The certificate in \cref{def:supp_normalized_numerical_certificate} is stable
under small smooth changes of the finite-response port.  More precisely, on
any compact window where
\[
  \mathcal A\le \frac9{20},\qquad
  \mathcal R\ge \frac94,\qquad
  \mathcal S\ge 4,\qquad
  \delta_{\rm sm}\le\frac{b_{\rm sep}}{24},\qquad
  \tau_{\rm strip}\le 8\cdot10^{-4}\delta_x,
\]
the smoothed two-mode port still satisfies the nonconvexity certificate after
the response-corrector, smoothing, and high-strip losses are subtracted.
\end{proposition}

\begin{proof}
The amplitude inequality keeps a strict geometric buffer
\(1-\mathcal A\ge11/20\), so the deformation remains inside the finite-horizon
window and the bump supports remain separated after smoothing.  The
stabilizer inequality gives the compact cotangent confinement used in
\cref{cor:supp_concrete_parameter_recipe}.

For the nonconvex Taylor coefficient, write the direct contribution as
\[
  2\lambda\langle\eta\rangle_{\theta_0}
  =\mathcal R(C_{\rm resp}+2\delta_x).
\]
After subtracting the response-corrector budget, the smoothing loss, and the
high-strip tail, the remaining gap is bounded below by
\[
  (\mathcal R-1)C_{\rm resp}+(2\mathcal R-1)\delta_x-\tau_{\rm strip}.
\]
With \(\mathcal R\ge9/4\) and
\(\tau_{\rm strip}\le 8\cdot10^{-4}\delta_x\), this lower bound is positive
and in particular exceeds \(\delta_x\) on the normalized window.  Hence the
directional second coefficient stays negative.  The inequalities are open in
the port coefficients because all read-outs and imported response constants
are continuous on the compact response window, so nearby smooth ports inherit
the same certificate.
\end{proof}

\begin{lemma}[Directional nonconvexity certificate]
\label{lem:supp_directional_nonconvexity}
Let \(e\) be a unit direction and suppose the directional Taylor coefficient
\[
  \mathfrak c_e(x)
  =
  \left.\frac{d^2}{ds^2}H(x,se)\right|_{s=0}
\]
is negative on an open set.  Then \(p\mapsto H(x,p)\) is locally nonconvex
there.  In the concrete two-mode port of the main paper it is enough that
\[
  2\lambda(x)\langle\eta\rangle_{\theta_0}
  >
  C_{\rm resp}(x)+\delta_x .
\]
\end{lemma}

\begin{proof}
Taylor expansion gives
\[
  H(x,se)+H(x,-se)-2H(x,0)
  =
  \mathfrak c_e(x)s^2+o(s^2),
\]
which is negative for small nonzero \(s\).  For the concrete port, the direct
action contribution to \(\mathfrak c_e\) is
\(-2\lambda(x)\langle\eta\rangle_{\theta_0}\).  The remaining invariant-density
and corrector-response terms are bounded by \(C_{\rm resp}(x)\).  The
dominance inequality therefore forces \(\mathfrak c_e(x)<0\).  By
\cref{prop:supp_concrete_parameter_window}, this is realized on an open set of
finite-response port coefficients before the transfer-resolvent calculation is
performed.
\end{proof}

\begin{lemma}[Generator-level subadditivity obstruction]
\label{lem:supp_subadditivity_obstruction}
Let \(S_t\phi=\Etheta_t[\phi]\).  If \(S_t\) were subadditive for all smooth
localized terminal payoffs, then the Hamiltonian would satisfy
\[
  H(x,p+q)\ge H(x,p)+H(x,q)
\]
with the terminal-time sign convention of the main article.  Hence any strict
opposite short-time Hamiltonian defect gives a failure of subadditivity.
\end{lemma}

\begin{proof}
For a smooth payoff with gradient \(p\) and zero Hessian at \(x\), the
terminal-value HJB gives
\[
  S_t\phi(x)=\phi(x)-tH(x,p)+o(t).
\]
Apply this to two localized affine payoffs with gradients \(p\) and \(q\), and
to their sum.  Subadditivity
\(S_t(\phi+\psi)\le S_t\phi+S_t\psi\) implies the displayed inequality after
dividing by \(t\) and passing \(t\downarrow0\).  Therefore a strict reverse
defect of the Hamiltonian at the chosen gradients produces the claimed failure.
\end{proof}

\begin{proposition}[Appendix-to-main crosswalk]
\label{prop:supp_hjb_crosswalk}
The technical blocks above support the following main-paper assertions.
\begin{center}
\small
\begin{tabular}{p{0.32\linewidth}p{0.58\linewidth}}
\hline
Main-paper assertion & Appendix support \\
\hline
Finite-response coefficients are microscopic &
\cref{def:supp_micro_action_port,prop:supp_no_effective_feedback} \\
Corrector hierarchy and residual identity &
\cref{def:supp_full_corrector_hierarchy,prop:supp_corrector_hierarchy_closure,lem:supp_first_order_cancellation,prop:supp_second_cell_average,lem:supp_residual_identity,def:supp_branch_residual_family,prop:supp_full_branch_residual_check,prop:supp_four_scale_ledger} \\
Half-relaxed HJB passage &
\cref{lem:supp_admissible_density_smoothing,prop:supp_perturbed_test_smoothing,cor:supp_smooth_approximation_closure,lem:supp_contact_selection,prop:supp_subsuper_passage,lem:supp_synchronized_density_doubling,prop:supp_doubled_variable_comparison,thm:supp_detailed_hjb_closure} \\
Coefficient comparison and differentiability &
\cref{prop:supp_coefficient_differentiability,def:supp_correlation_hilbert_bundle,prop:supp_correlation_factor_transport,prop:supp_gk_comparison_trace,prop:supp_comparison_modulus_ledger,prop:supp_terminal_stability} \\
Nonconvexity and non-subadditivity &
\cref{prop:supp_concrete_parameter_window,cor:supp_concrete_parameter_recipe,def:supp_dimensionless_parameter_audit,prop:supp_parameter_audit_verification,def:supp_normalized_numerical_certificate,prop:supp_numerical_certificate_margin,lem:supp_directional_nonconvexity,lem:supp_subadditivity_obstruction} \\
\hline
\end{tabular}
\end{center}
\end{proposition}

